\begin{table}[!tbp]
\caption{\textbf{Core ablation of progressive static-to-dynamic reasoning.} The upper block compares prediction paradigms; the lower block isolates suppression and raw-feature bypass. The parallel S2D control uses the same no-static-suppression setting as the lower-block control.}
\label{tab:core_mechanism}
\centering
\scriptsize
\setlength{\tabcolsep}{4.2pt}
\renewcommand{\arraystretch}{1.08}
\begin{tabularx}{\columnwidth}{@{}>{\raggedright\arraybackslash}Xcccc@{}}
\toprule
Setting & gIoU$\uparrow$ & $\mathrm{mIoU}_{\mathrm{all}}\uparrow$ & $\mathrm{mIoU}_{\mathrm{dyn}}\uparrow$ & $\mathrm{mIoU}_{\mathrm{sta}}\uparrow$ \\
\midrule
\multicolumn{5}{l}{\textbf{\textit{Necessity of Progressive S2D Reasoning}}} \\
Plain Occupancy Prediction & 78.54 & 50.71 & 23.48 & 71.14 \\
Parallel S2D Fusion (no supp.) & 79.77 & 53.59 & 24.43 & 75.46 \\
Progressive S2D Reasoning & 89.08 & 60.08 & 28.97 & 83.42 \\
\midrule
\multicolumn{5}{l}{\textbf{\textit{Design of Progressive S2D Reasoning}}} \\
No static suppression & 79.77 & 53.59 & 24.43 & 75.46 \\
Suppression only & 88.35 & 59.77 & 28.81 & 82.99 \\
Suppression + raw bypass & 89.08 & 60.08 & 28.97 & 83.42 \\
\bottomrule
\end{tabularx}
\end{table}
